\documentclass[11pt]{article}

\usepackage[margin=1in]{geometry}
\usepackage{amsmath,amssymb,amsthm,mathtools,booktabs,tabularx}
\usepackage{microtype,xurl}
\usepackage{xcolor}
\usepackage[colorlinks=true,linkcolor=blue!55!black,
            citecolor=blue!55!black,urlcolor=blue!55!black]{hyperref}

\newtheorem{theorem}{Theorem}[section]
\newtheorem{proposition}[theorem]{Proposition}
\newtheorem{corollary}[theorem]{Corollary}
\theoremstyle{definition}
\newtheorem{definition}[theorem]{Definition}
\theoremstyle{remark}
\newtheorem{remark}[theorem]{Remark}

\newcommand{\dd}{\mathrm d}
\newcommand{\Lie}{\mathcal L}
\newcommand{\Diff}{\operatorname{Diff}}
\newcommand{\Weyl}{\operatorname{Weyl}}
\newcommand{\cert}[1]{\nolinkurl{#1}}
\newcolumntype{Y}{>{\raggedright\arraybackslash}X}

\title{Residual-Basic Charges and Simultaneous Horizon First Laws\\
\large in Static Weyl Gravity}
\author{GPT-5.6.sol\thanks{OpenAI model and principal programme author.}
\and Asger Alstrup Palm\thanks{Programme orchestrator and corresponding human contact; Honorary Professor, DTU Compute, Technical University of Denmark. \texttt{asger@area9.dk}.}}
\date{Working draft, 26 July 2026}

\begin{document}
\maketitle

\begin{abstract}
Static spherical solutions of four-dimensional Weyl gravity possess a
residual local \(\Diff\times\Weyl\) freedom that changes the normalization of
the chart Killing vector.  We show, by exact covariant-phase-space
calculation, that this freedom resolves an apparent nonintegrability of the
static Iyer--Wald charge.  On the Mannheim--Kazanas component of a
Laurent-complete Bach-flat family, the bare charge associated with
\(\partial_t\) is not closed.  Requiring the field-dependent generator and
its charge to descend to the residual quotient forces
\[
 \chi=u f(J)\partial_t,\qquad
 u=\beta(2-3\beta\gamma),
\]
where \(J\) is the single continuous residual invariant.  For the canonical
choice \(f=1\), the corrected charge is exact, with Hamiltonian
\[
 H=-16\pi\alpha\beta^2
 \bigl(9\beta^2\gamma^2k-\beta\gamma^3-12\beta\gamma k
       +\gamma^2+4k\bigr).
\]
At every simple root \(r_h\) of the metric function, the Wald entropy and
signed temperature
\[
 S_h=64\pi^2\alpha\beta
 \frac{2-3\beta\gamma+\gamma r_h}{r_h},
 \qquad
 T_h=\frac{uB'(r_h)}{4\pi}
\]
obey \(\dd H=T_h\dd S_h\).  The same one-form \(\dd H\) therefore gives
simultaneous first laws at all simple horizons, including an exact
non-Einstein three-horizon example.  At linear charge level the result is
unchanged by arbitrary time-dependent spherical Weyl rescalings and
diffeomorphisms: those directions carry zero charge, zero entropy variation,
and zero corrected presymplectic pairing with the static parameter modes.

The theorem is deliberately scoped.  The background classification is
complete only in the declared Laurent class, the charge theorem is exact on
the static parameter family and in the linear spherical gauge sector, and
no preferred physical mass, second-order physical-process first law,
radiative flux law, stability statement, or quantum claim is inferred.
\end{abstract}

\section*{AI authorship and computational accountability}

GPT-5.6.sol developed the derivations, symbolic implementations, verification
architecture, claim audit, and manuscript.  Asger Alstrup Palm commissioned
and orchestrated the programme and is the corresponding human contact.  The
project is also an experiment in whether a non-specialist human, using AI
systems as research collaborators, can organize falsifiable
computer-assisted mathematical physics.  Every promoted identity in this
paper is bound to a machine-readable exact certificate and to a separately
implemented verifier.  The human orchestrator has not independently checked
every tensor calculation; the manuscript therefore remains subject to
expert review.

\tableofcontents

\section{Question, contribution, and claim boundary}

For a stationary black hole in an asymptotically flat theory, the
normalization of the horizon generator is usually fixed by the asymptotic
clock.  Generic static spherical solutions of pure Weyl gravity are not
asymptotically flat, and their Schwarzschild-gauge representatives retain a
local radial-diffeomorphism/Weyl freedom.  The chart vector
\(\partial_t\) is consequently not a residual-basic object.

This matters already on the three-dimensional static parameter family.  The
Iyer--Wald surface form for \(\partial_t\) is radially conserved but not
integrable in parameter space.  Adding a parameter-local corner function
cannot cure this, because it cannot change the exterior derivative of the
charge one-form.  Our central observation is that the obstruction disappears
when the generator is normalized as an object on the residual quotient.
Basicness determines the normalization up to a function of the sole quotient
invariant, and the resulting charge is closed.

The static Bach-flat family itself is classical
\cite{Riegert,MannheimKazanas}; conformal-gravity black-hole thermodynamics
and additional spin-two hair have also been studied in asymptotically AdS
and Einstein--Weyl settings \cite{LuPangPopeVP}.  The contribution here is
not a new solution family.  It is:
\begin{enumerate}
\item an exact residual-quotient derivation of the field-dependent
      normalization;
\item an exact Hamiltonian and Wald entropy in that normalization;
\item one first-law identity holding simultaneously at every simple horizon;
\item a linear spherical audit showing that the charge is not selected by a
      conformal frame or by a spherical coordinate gauge.
\end{enumerate}
L\"u, Pang, Pope, and V\'azquez-Poritz use an AdS ensemble with an
independent thermodynamic pair for massive spin-two hair
\cite{LuPangPopeVP}.  Our one-term identity is instead a quotient statement
on a restricted static family with no preferred asymptotic clock.  The two
constructions therefore answer different normalization and ensemble
questions.

\subsection{Dependency closure}

Table~\ref{tab:dependencies} separates the exact theorem chain from open
extensions.  The dependency tags are those used in the computational
archive.  In the compact status column, LA means
\textsc{local-algebraic} and RM means \textsc{reduced-mode}; the abbreviated
certificate names are expanded in Section~9.

\begin{table}[ht]
\centering
\small
\begin{tabularx}{\textwidth}{@{}p{0.29\textwidth}p{0.10\textwidth}p{0.18\textwidth}Y@{}}
\toprule
Result & Evidence & Status & Role in main theorem\\
\midrule
Laurent Bach-flat locus, residual action, horizon fixture
& BH0
& LA, exact
& Background and quotient geometry\\
Bare static Iyer--Wald form
& BH1
& LA+RM, exact
& Exhibits the normalization obstruction\\
Basic generator, \(H,S,T\), simultaneous first law
& BH1A
& LA+RM, exact
& Main static theorem\\
Spherical conformal/diffeomorphism audit
& BH1B
& LA+RM, exact at linear charge level
& Frame and gauge independence\\
Radiative flux, physical-process law, stability
& none
& open
& Not used and not claimed\\
\bottomrule
\end{tabularx}
\caption{Logical dependency of the paper.  ``Exact'' means symbolic rational
or polynomial identity, not floating-point evidence.}
\label{tab:dependencies}
\end{table}

\section{Conventions and the static Laurent family}

We use signature \((-+++)\) and action
\begin{equation}\label{eq:action}
 S_W[g]=\alpha\int_M\sqrt{-g}\,
 C_{abcd}C^{abcd}\,\dd^4x.
\end{equation}
Our Bach-tensor convention is
\begin{equation}
 B_{ab}=\nabla^c\nabla^dC_{acbd}
       +\frac12R^{cd}C_{acbd}.
\end{equation}
The field equation is \(B_{ab}=0\).

Start from the areal-radius ansatz
\begin{equation}
 \dd s^2=-a(r)\dd t^2+b(r)\dd r^2+r^2\dd\Omega_2^2.
\end{equation}
Before conformal gauge fixing, the exact Bach tensor is diagonal,
trace-free, divergence-free, and has two independent components.  We then
choose \(b=1/a\), write \(a=B\), and restrict to
\begin{equation}\label{eq:laurent}
 B(r)=w-\frac{u}{r}+\gamma r-kr^2
      +\frac{c_2}{r^2}+c_3r^3.
\end{equation}

\begin{theorem}[Bach-flat locus in the declared Laurent class]
\label{thm:laurent}
The metric \eqref{eq:laurent} is Bach-flat if and only if
\begin{equation}\label{eq:laurent-locus}
 c_2=c_3=0,\qquad w^2+3u\gamma=1.
\end{equation}
\end{theorem}

\begin{proof}
Substitution of \eqref{eq:laurent} into the two independent Bach rows,
clearing the declared denominators and equating Laurent coefficients, gives
an exact polynomial ideal.  Its reduced Gr\"obner basis is
\[
 \{c_2,\ c_3,\ w^2+3u\gamma-1\}.
\]
This proves necessity and sufficiency within the ansatz.  The producer and
independent curvature implementation both recompute the vanishing tensor;
mutating the correlated constant term or adding an \(r^{-2}\) tail produces
explicit nonzero Bach components.
\end{proof}

\begin{remark}[Scope of completeness]
Theorem~\ref{thm:laurent} is not presented as a new proof of the full
Birkhoff--Riegert theorem \cite{Riegert}.  Generalized Birkhoff statements
and nonstandard spherical branches require additional global and
nondegeneracy hypotheses \cite{JizbaLeheckova}.  Here ``complete'' always
means complete in \eqref{eq:laurent}.
\end{remark}

The Mannheim--Kazanas chart is
\begin{equation}\label{eq:mk}
 B(r)=1-3\beta\gamma
      -\frac{\beta(2-3\beta\gamma)}{r}
      +\gamma r-kr^2,
\end{equation}
so
\begin{equation}\label{eq:uw}
 u=\beta(2-3\beta\gamma),\qquad
 w=1-3\beta\gamma.
\end{equation}
It parametrizes the component through \(w=1\) used in the charge theorem.

\begin{proposition}[Correct Einstein locus]
\label{prop:einstein}
On the full Laurent locus \eqref{eq:laurent-locus}, the metric is Einstein
in the areal-radius Schwarzschild chart if and only if
\begin{equation}
 \gamma=0,\qquad w=1.
\end{equation}
On the Mannheim--Kazanas component \eqref{eq:mk}, this reduces to
\(\gamma=0\).
\end{proposition}

\begin{proof}
On the full locus the scalar curvature is
\begin{equation}
 R=12k-\frac{6\gamma}{r}+\frac{2(1-w)}{r^2}.
\end{equation}
An Einstein metric has constant scalar curvature, forcing
\(\gamma=0\) and \(w=1\).  Direct substitution then gives
\(R_{ab}=3k g_{ab}\).  In the MK chart, the trace-free Ricci defect contains
\[
 E_{\theta\theta}=-\frac{\gamma}{2}(r-3\beta),
\]
so \(\gamma=0\) is necessary and sufficient there.  The allowed Laurent
point \((\gamma,w)=(0,-1)\) is therefore Bach-flat but not Einstein; it
belongs to a complementary Kantowski--Sachs-type chart rather than to the
MK component through \(w=1\).
\end{proof}

\section{Residual quotient and its invariant}

Put \(x=1/r\) and
\begin{equation}
 Q(x)=-ux^3+wx^2+\gamma x-k.
\end{equation}
The roots of \(Q\) encode the nonzero roots of \(B\).  The form-preserving
local transformation
\begin{equation}
 r=\frac{\rho}{1-c\rho},\qquad
 g\longmapsto(1-c\rho)^2g
\end{equation}
acts as the translation \(Q(x)\mapsto Q(x-c)\).  Equivalently,
\begin{align}
 u&\longmapsto u,&
 w&\longmapsto w+3cu,\nonumber\\
 \gamma&\longmapsto\gamma-2cw-3c^2u,&
 k&\longmapsto k+c\gamma-c^2w-c^3u.
\label{eq:c-action}
\end{align}
A dilation acts by
\begin{equation}
 (\beta,\gamma,k)\longmapsto
 (\beta/\lambda,\lambda\gamma,\lambda^2k),
 \qquad
 Q(x)\longmapsto\lambda^2Q(x/\lambda).
\end{equation}
The conformal factor in the \(c\)-map may vanish at finite radius, so
\eqref{eq:c-action} is used as a local quotient statement, not as a global
identification of black-hole exteriors.

On the MK parameters the infinitesimal generators are
\begin{align}
 X_c&=-3\beta^2\partial_\beta
 +(6\beta\gamma-2)\partial_\gamma+\gamma\partial_k,
\label{eq:Xc}\\
 X_\lambda&=-\beta\partial_\beta
 +\gamma\partial_\gamma+2k\partial_k.
\label{eq:Xl}
\end{align}
They have generic rank two.  A continuous invariant is
\begin{equation}\label{eq:J}
 J=u^2\operatorname{disc}(Q).
\end{equation}
Translation invariance of the discriminant and its scaling weight prove
\(X_cJ=X_\lambda J=0\).  Moreover,
\(\operatorname{disc}(Q)=0\) is precisely the multiple-root, hence
degenerate-horizon, locus.  Locally on a regular orbit component, every
continuous invariant is therefore a function of \(J\).

\section{The bare charge is not integrable}

For a diffeomorphism-covariant Lagrangian depending algebraically on the
Riemann tensor, define
\[
 E^{abcd}=\frac{\partial(\alpha C_{efgh}C^{efgh})}
                 {\partial R_{abcd}}
          =2\alpha C^{abcd}.
\]
We use the Iyer--Wald representative \cite{LeeWald,IyerWald}
\begin{align}
 \theta^a(g;\delta g)
 &=2\left(E^{abcd}\nabla_d\delta g_{bc}
          -\delta g_{bc}\nabla_dE^{abcd}\right),\\
 Q^{ab}_\xi
 &=-E^{abcd}\nabla_c\xi_d+2\xi_d\nabla_cE^{abcd}.
\end{align}
The sphere conventions are
\[
 (Q_\xi)_{\theta\phi}=2\sqrt{-g}\,Q^{tr}_\xi,\qquad
 (i_\xi\theta)_{\theta\phi}=-\sqrt{-g}\,\xi^t\theta^r.
\]
With \(\xi=\partial_t\), integration over any sphere gives the
parameter-space one-form
\begin{equation}\label{eq:F}
 \mathcal F=F_\beta\dd\beta+F_\gamma\dd\gamma+F_k\dd k,
\end{equation}
where
\begin{align}
 F_\beta&=16\pi\alpha(12\beta\gamma k-\gamma^2-4k),\nonumber\\
 F_\gamma&=16\pi\alpha\beta(6\beta k-\gamma),\label{eq:Fcomponents}\\
 F_k&=-16\pi\alpha\beta(2-3\beta\gamma).\nonumber
\end{align}
Each component is exactly independent of the sphere radius.  As a
normalization control, replacing \(E^{abcd}\) by the Einstein--Hilbert
tensor reproduces \(16\pi\,\delta m\) on Schwarzschild and
Schwarzschild--de Sitter in the convention \(16\pi G=1\).

\begin{proposition}[Bare nonintegrability]
\label{prop:bare}
The chart-normalized form \(\mathcal F\) is not closed:
\begin{equation}\label{eq:dF}
 \dd\mathcal F=16\pi\alpha\left[
 \gamma\,\dd\beta\wedge\dd\gamma
 +2(1-3\beta\gamma)\,\dd\beta\wedge\dd k
 -3\beta^2\,\dd\gamma\wedge\dd k\right].
\end{equation}
It obeys
\begin{align}
 i_{X_c}\mathcal F&=0,&i_{X_c}\dd\mathcal F&=0,\label{eq:c-horizontal}\\
 i_{X_\lambda}\mathcal F&=0,&
 i_{X_\lambda}\dd\mathcal F&=\mathcal F.\label{eq:l-euler}
\end{align}
Consequently no parameter-local boundary function can make
\(\mathcal F\) integrable.
\end{proposition}

\begin{proof}
Equations \eqref{eq:dF}--\eqref{eq:l-euler} follow by exact
differentiation of \eqref{eq:Fcomponents}.  If a boundary function \(W\)
were added, the one-form would change by \(\dd W\), leaving
\(\dd\mathcal F\) unchanged.  The kernel of \(\dd\mathcal F\) is the
\(X_c\) direction at generic points.  Thus quotienting only the
translation leaves a nondegenerate two-form on the remaining
two-dimensional parameter quotient; the obstruction is not removed merely
by discarding that proper-gauge direction.
\end{proof}

\section{Basicness fixes the field-dependent generator}

For a field-dependent generator \(\chi=N(\beta,\gamma,k)\partial_t\), we
use the corrected surface form
\begin{equation}\label{eq:fdcharge}
 k_\chi(g;\delta g)
 =\delta Q_\chi-Q_{\delta\chi}-i_\chi\theta(g;\delta g).
\end{equation}
Because \(N\) is spacetime-constant on each static solution and the Noether
charge is linear in \(\chi\), the \(Q_{\delta\chi}\) term cancels all
\(\dd N\) contributions.  The corrected charge is therefore
\begin{equation}
 \delta H_\chi=N\mathcal F.
\end{equation}

\begin{definition}
A parameter-space form is \emph{residual-basic} if it is horizontal and
invariant under both vector fields \eqref{eq:Xc} and \eqref{eq:Xl}.
\end{definition}

\begin{theorem}[Residual-basic normalization]
\label{thm:normalization}
On a connected regular component with \(u\neq0\), the corrected charge
\(N\mathcal F\) is residual-basic only if
\begin{equation}\label{eq:N}
 N=u f(J)
\end{equation}
for a function \(f\) of the quotient invariant.  Conversely, every
normalization \eqref{eq:N} is basic and integrable.  For \(f=1\),
\begin{equation}\label{eq:H}
 u\mathcal F=\dd H,\qquad
 H=-16\pi\alpha\beta^2D_2,
\end{equation}
where
\begin{equation}\label{eq:D2}
 D_2=9\beta^2\gamma^2k-\beta\gamma^3
     -12\beta\gamma k+\gamma^2+4k.
\end{equation}
\end{theorem}

\begin{proof}
Cartan's identity and \eqref{eq:c-horizontal} give
\(\Lie_{X_c}\mathcal F=0\).  Hence invariance of \(N\mathcal F\) requires
\(X_cN=0\).  Similarly, \eqref{eq:l-euler} gives
\(\Lie_{X_\lambda}\mathcal F=\mathcal F\), so invariance requires
\(X_\lambda N=-N\).  Directly,
\[
 X_cu=0,\qquad X_\lambda u=-u.
\]
Since the residual orbit has codimension one and \(J\) is its local
invariant, the general local solution is \(N=uf(J)\).

Exact differentiation of \eqref{eq:H} gives
\(\dd H=u\mathcal F\), proving closure for \(f=1\).  The same form is
horizontal under both residual fields.  For general \(f(J)\),
\(f(J)\dd H\) is closed because the exact identity
\(\dd H\wedge\dd J=0\) holds.  It is therefore locally the differential of
a reparametrized quotient Hamiltonian.
\end{proof}

The discriminant factorization sharpens the quotient interpretation.  Put
\begin{equation}
 D_1=27\beta^2k-3\beta\gamma-1.
\end{equation}
Then
\begin{equation}\label{eq:Jfactor}
 J=-u^2D_1D_2.
\end{equation}
The Hamiltonian vanishes on the \(D_2=0\) component of the degenerate-horizon
locus.  On the Einstein control family \(\gamma=0\),
\begin{equation}
 H=-64\pi\alpha\beta^2k.
\end{equation}
On the Schwarzschild subfamily \(k=\gamma=0\), it vanishes identically.

\begin{remark}[Normalization versus physical mass]
Equation \eqref{eq:H} defines the Hamiltonian of the residual-basic
generator on the declared static phase-space slice.  We do not call it a
universal physical mass.  Multiplication by \(f(J)\), component orientation,
and any future choice of physical asymptotic clock change the normalization.
On a fixed-falloff ensemble with \(\gamma\) and \(k\) fixed, the nontrivial
residual transformations are generically frozen, but this does not create a
preferred clock on the full family.
\end{remark}

\section{Entropy and simultaneous first laws}

Let \(r_h\) be any simple positive root of \(B\).  The Wald entropy is
\begin{equation}\label{eq:wald}
 S_h=-2\pi\int_{S_{r_h}}E^{abcd}\epsilon_{ab}\epsilon_{cd}\,\dd A.
\end{equation}
With \(E^{abcd}=2\alpha C^{abcd}\) and \(\epsilon_{tr}=1\), exact evaluation
gives
\begin{equation}\label{eq:S}
 S_h=64\pi^2\alpha\beta
 \frac{2-3\beta\gamma+\gamma r_h}{r_h}.
\end{equation}
For the algebraic orientation \(\chi=u\partial_t\), define the signed
surface-gravity temperature
\begin{equation}\label{eq:T}
 T_h=\frac{\kappa_\chi}{2\pi}
     =\frac{uB'(r_h)}{4\pi}.
\end{equation}
This is the signed quantity entering the oriented first law; a positive
Hawking temperature would use \(|uB'(r_h)|/(4\pi)\).

\begin{theorem}[Simultaneous horizon first laws]
\label{thm:firstlaw}
For every simple root \(r_h\) of \eqref{eq:mk}, the functions
\eqref{eq:H}, \eqref{eq:S}, and \eqref{eq:T} satisfy
\begin{equation}\label{eq:firstlaw}
 \boxed{\dd H=T_h\,\dd S_h}
\end{equation}
in all three parameter directions.  If one solution has simple roots
\(r_1,\ldots,r_n\), then
\begin{equation}\label{eq:simultaneous}
 T_1\dd S_1=\cdots=T_n\dd S_n=\dd H.
\end{equation}
\end{theorem}

\begin{proof}
Along the horizon locus \(B(r_h)=0\), any parameter \(p\in
\{\beta,\gamma,k\}\) obeys
\begin{equation}
 \frac{\dd r_h}{\dd p}=-\frac{\partial_pB(r_h)}{B'(r_h)}.
\end{equation}
Use this in the total derivative of \eqref{eq:S}, subtract it from
\(\partial_pH/T_h\), clear denominators, and reduce the numerator modulo
\begin{equation}
 r_hB(r_h)=0.
\end{equation}
The remainder is identically zero for each of the three parameters.  This
calculation uses no choice of root beyond simplicity, proving the result at
all roots simultaneously.
\end{proof}

For Schwarzschild, \(r_h=2\beta\), so
\[
 S_h=64\pi^2\alpha
\]
is independent of \(\beta\), consistently with \(H=0\) on that
subensemble.  No pressure--volume or extra source term is needed for
\eqref{eq:firstlaw} on the declared residual-basic static family.  This
does not exclude such terms in a different ensemble with independently
varied boundary data.

\subsection{Exact non-Einstein three-horizon fixture}

Consider
\begin{equation}\label{eq:fixture}
 (\beta,\gamma,k)=\left(\frac32,\frac{12}{19},\frac1{19}\right).
\end{equation}
Then
\begin{equation}
 rB(r)=-\frac1{19}(r-1)(r-3)(r-8),\qquad
 u=-\frac{24}{19},
\end{equation}
and the trace-free Ricci witness is
\[
 E_{\theta\theta}=-\frac{3(2r-9)}{19}\neq0.
\]
The metric is smooth at all three roots in ingoing
Eddington--Finkelstein coordinates, and all certified scalar curvature
invariants are finite there.  In the algebraic \(N=u\) orientation,
\begin{equation}
 H=-\frac{14400\pi\alpha}{6859}
\end{equation}
and
\begin{center}
\begin{tabular}{c|cc}
\toprule
\(r_h\) & \(S_h\) & \(T_h\)\\
\midrule
\(1\) & \(-384\pi^2\alpha/19\) & \(84/(361\pi)\)\\
\(3\) & \(640\pi^2\alpha/19\) & \(-20/(361\pi)\)\\
\(8\) & \(960\pi^2\alpha/19\) & \(105/(1444\pi)\)\\
\bottomrule
\end{tabular}
\end{center}
The three identities \(T_h\dd S_h=\dd H\) hold exactly in every parameter
direction.  Since \(u<0\) on this component, choosing the opposite
future-directed orientation flips \(H\) and every signed \(T_h\), leaving
the first laws unchanged.

\section{Linear spherical frame and gauge audit}

The preceding derivation lives on the static parameter slice.  It is
nevertheless important to decide whether its normalization secretly
selects a conformal frame.  We therefore consider the complete linear
spherical gauge sector around both symbolic Schwarzschild and the
non-Einstein fixture:
\begin{align}
 \delta_\omega g_{ab}&=2\omega(t,r)g_{ab},\\
 \delta_\xi g_{ab}&=\Lie_\xi g_{ab},\qquad
 \xi=a(t,r)\partial_t+b(t,r)\partial_r.
\end{align}

\begin{theorem}[Linear spherical gauge independence]
\label{thm:l0}
At linear charge level:
\begin{enumerate}
\item \(k_\chi(2\omega g)=0\) componentwise for arbitrary \(\omega(t,r)\);
\item the entropy variation under \(2\omega g\) vanishes on the symbolic
      MK family;
\item the corrected presymplectic current between a conformal variation and
      any static parameter mode vanishes;
\item \(k_\chi(\Lie_\xi g)=0\) for arbitrary spherical \(\xi\);
\item the extension of \(N=u\) is uniquely fixed by
      \(\delta N=\dd u\) on parameter directions and \(\delta N=0\) on
      spherical conformal and diffeomorphism directions.
\end{enumerate}
Thus no boundary clock or preferred conformal representative is needed for
the static theorem at this linear charge level.
\end{theorem}

\begin{proof}
The conformal statements follow from direct substitution in
\eqref{eq:fdcharge}, exact conformal covariance of the Bach equation, and
the literal variation of the density \(\sqrt{-g}\theta^a\).  The density
variation is essential: omitting the trace contribution spoils current
conservation.  The diffeomorphism statement is verified through the
on-shell Noether identity
\begin{equation}
 \theta(\Lie_\xi g)-i_\xi(L\epsilon)=\dd Q_\xi
\end{equation}
and the background identity obtained from its commutator with \(\chi\).
Both routes are independently evaluated in the certificate.

For uniqueness, the bare Noether aperture is nonzero on the fixture.  Any
nonzero \(\delta N\) on a certified gauge direction would therefore assign
that direction a nonzero charge, contradicting the componentwise
annihilation above.  The parameter variation of \(N\) is already fixed by
Theorem~\ref{thm:normalization}.
\end{proof}

The theorem is complete for the \(\ell=0\) dynamical gauge sector at linear
charge level.  It does not compute the bilinear radiative flux for
\(\ell\ge2\), and it does not imply a second-order physical-process first
law.

\section{Ensembles, interpretation, and non-results}

The \(c\)-transformation preserves the fixed-falloff data
\(\{\gamma,k\}\) only for \(c=0\), away from the algebraic locus
\[
 u(w^2-3)=0.
\]
The \(u=0\) component is excluded because the basic normalization
degenerates.  The dilation preserves the same ensemble only at
\(\lambda=1\), except on the Schwarzschild subensemble
\(\gamma=k=0\), where it acts freely and consistently with \(H=0\) and
constant entropy.  On a finite exterior \([r_h,r_{\rm out}]\), the
\(c\)-map conformal factor is smooth and positive precisely when
\(1+cr>0\) throughout that interval.

These observations support a precise interpretation:
\begin{quote}
The nonintegrability of the chart-normalized charge is a normalization
artifact.  Requiring the generator and charge to descend to the complete
local residual quotient selects an integrable class of generators.  It does
not, by itself, choose a unique physical clock.
\end{quote}

The following statements are not established:
\begin{itemize}
\item completeness of static Bach vacua beyond the Laurent ansatz;
\item a preferred physical matter or clock conformal frame;
\item a universal physical mass valid across inequivalent boundary
      ensembles;
\item a nonlinear or second-order physical-process first law;
\item the bilinear radiative symplectic flux through a horizon or infinity;
\item stability, quasinormal ringing, Hawking radiation, or any quantum
      property.
\end{itemize}

\section{Computational evidence and reproduction}

The calculations are computer-assisted but exact.  Tensor components,
polynomial reductions, and charge identities are represented by symbolic
rational functions.  No theorem in this paper is inferred from a decimal
approximation.  Producer and verifier use separately written curvature
pipelines; the verifier constructs the Weyl tensor through an independent
Schouten/Kulkarni--Nomizu route.  Mutation tests ensure that certificate-like
but false inputs are rejected.

\begin{table}[ht]
\centering
\small
\begin{tabularx}{\textwidth}{@{}p{0.22\textwidth}Y@{}}
\toprule
Certificate & Reproduction and independent verification\\
\midrule
\cert{BH0_STATIC_SPHERICAL_BACKGROUND.json}
& \texttt{python3 black\_hole\_programme/bh0\_background.py}\\
& \texttt{python3 black\_hole\_programme/verify\_bh0\_background.py}\\[1mm]
\cert{BH1_LEE_WALD_PREFLIGHT.json}
& \texttt{python3 black\_hole\_programme/bh1\_lee\_wald\_preflight.py}\\
& \texttt{python3 black\_hole\_programme/verify\_bh1\_lee\_wald\_preflight.py}\\[1mm]
\cert{BH1A_NORMALIZED_GENERATOR.json}
& \texttt{python3 black\_hole\_programme/bh1a\_normalized\_generator.py}\\
& \texttt{python3 black\_hole\_programme/verify\_bh1a\_normalized\_generator.py}\\[1mm]
\cert{BH1B_DYNAMICAL_EXTENSION.json}
& \texttt{python3 black\_hole\_programme/bh1b\_dynamical.py}\\
& \texttt{python3 black\_hole\_programme/verify\_bh1b\_dynamical.py}\\
\bottomrule
\end{tabularx}
\caption{Authoritative certificate families.  Re-running a producer is
reproduction; the separately implemented verifier is the independent
check.}
\end{table}

The BH0 mutations (i) remove the correlated \(-3\beta\gamma\) constant and
(ii) add an \(r^{-2}\) tail to the three-horizon fixture; both yield nonzero
Bach components.  The bare-charge computation includes Einstein
normalization controls and retains \(\dd\mathcal F\neq0\) as a negative
control.  The BH1A verifier independently recomputes entropy and reduces the
first-law polynomial at all three fixture roots.  The BH1B verifier
recomputes the conformal and diffeomorphism annihilation identities.

\section{Conclusion}

Static Weyl gravity does not supply a preferred Killing normalization merely
by presenting a Schwarzschild-gauge metric.  On the certified
Mannheim--Kazanas family, treating \(\partial_t\) as fixed produces a
radially conserved but nonintegrable Iyer--Wald form.  Treating the generator
as field dependent and demanding residual basicness instead forces
\(\chi=u f(J)\partial_t\).  Its charge is integrable, its Wald entropy is
explicit, and one Hamiltonian variation satisfies a signed first law at
every simple horizon simultaneously.

The result is stronger than a coordinate identity and weaker than a complete
black-hole thermodynamics.  It is invariant under the complete linear
spherical conformal and diffeomorphism gauge sector, but it does not fix a
physical asymptotic clock or control nonlinear radiative flux.  Those
boundaries are part of the theorem, not deferred qualifications.

\begingroup
\raggedright
\begin{thebibliography}{99}

\bibitem{Bach}
R.~Bach,
\emph{Zur Weylschen Relativit\"atstheorie und der Weylschen Erweiterung des
Kr\"ummungstensorbegriffs},
Math. Z. \textbf{9} (1921) 110--135.

\bibitem{Riegert}
R.~J.~Riegert,
\emph{Birkhoff's theorem in conformal gravity},
Phys. Rev. Lett. \textbf{53} (1984) 315--318.

\bibitem{MannheimKazanas}
P.~D.~Mannheim and D.~Kazanas,
\emph{Exact vacuum solution to conformal Weyl gravity and galactic rotation
curves},
Astrophys. J. \textbf{342} (1989) 635--638.

\bibitem{LeeWald}
J.~Lee and R.~M.~Wald,
\emph{Local symmetries and constraints},
J. Math. Phys. \textbf{31} (1990) 725--743.

\bibitem{IyerWald}
V.~Iyer and R.~M.~Wald,
\emph{Some properties of Noether charge and a proposal for dynamical black
hole entropy},
Phys. Rev. D \textbf{50} (1994) 846--864,
\href{https://arxiv.org/abs/gr-qc/9403028}{arXiv:gr-qc/9403028}.

\bibitem{LuPangPopeVP}
H.~L\"u, Y.~Pang, C.~N.~Pope, and J.~F.~V\'azquez-Poritz,
\emph{AdS and Lifshitz black holes in conformal and Einstein--Weyl
gravities},
Phys. Rev. D \textbf{86} (2012) 044011,
\href{https://arxiv.org/abs/1204.1062}{arXiv:1204.1062}.

\bibitem{JizbaLeheckova}
P.~Jizba and T.~Lehe\v{c}kov\'a,
\emph{Generalized Birkhoff theorems and \(2+2\) direct product spacetimes in
Weyl conformal gravity},
Phys. Rev. D \textbf{113} (2026) 044060,
\href{https://doi.org/10.1103/9g8r-lrjk}{doi:10.1103/9g8r-lrjk},
\href{https://arxiv.org/abs/2512.18482}{arXiv:2512.18482}.

\end{thebibliography}
\endgroup

\end{document}
